\section{个人}
\begin{enumerate}
  \item 三法三制度。国籍法、出入境管理法和引渡法（外交保护、引渡和庇护）。
  \item 国籍制度。\zd{取得、丧失}。\bs{做题先看定居状况}。
        \begin{enumerate}
          \item 取得。出生、加入。
                \begin{enumerate}
                  \item 中国人的孩子，\zd{定居国外 + 已取得外国国籍除外}；父母\zd{无国籍 \hz 国籍不明 + 定居中国}。
                        \al{父母外国籍，定居中国，小孩有中国国籍 \cuo}。
                  \item 申请。不得再保留外国国籍。
                \end{enumerate}
          \item 丧失。混合原则。自动、非自动。
                \begin{enumerate}
                  \item \zd{自动丧失 = 定居外国 + 加入 \hz 取得外国国籍} \yjt 无需办理手续。
                  \item 依申请。\zd{国家工作人员和现役军人不得退出中国国籍}。
                \end{enumerate}
          \item \zd{公安部}对入籍和退籍审批。
          \item 华侨。指定居外国的中国公民。会注销户口，没有身份证，护照是其身份证明。
          \item 冲突的解决。\bs{不能直接以经居地代替国籍国}。
                \begin{enumerate}
                  \item 积极。
                        \begin{enumerate}
                          \item 有中国国籍 \yjt 中国，其他不认。
                          \item 没有，多个国籍中的经居地 \yjt 最密切联系。
                        \end{enumerate}
                  \item 消极。无国籍看\zd{经居地国}。
                \end{enumerate}
        \end{enumerate}
  \item 出入境管理制度。境外签证驻外使领馆、外交部委托的其他驻外机构。出入境检查边检机关。入境后\zd{县以上公安}（有处罚权，可以拘留、审查）。
        \begin{enumerate}
          \item 签证制度。
                \begin{enumerate}
                  \item \zd{4} 大类，\zd{外交、礼遇、公务、普通}签证，普通再分 12 类。
                  \item 免签 3 个等。\zd{免签协议、中国居留证、过境不越界不超时}等。
                  \item 拒签 4 个等。\zd{无需说明理由}。\al{精神障碍，不予签证 \dui}。
                \end{enumerate}
          \item 居留制度。\zd{打击 3 非。非法入境、非法居留、非法工作}。
                \begin{enumerate}
                  \item \zd{邀请函件真实性}，\zd{单位和个人}也要承担法律责任。
                  \item 住宿登记 \zd{24} 小时。
                  \item 工作类签证、居留证。\bdy 公务签证。例外。\zd{留学生}实习、勤工助学。学校同意居留证件加注相关信息。
                \end{enumerate}
          \item 出境。边检有权阻止出境。
                \begin{enumerate}
                  \item \zd{限制出境 3 个等}。涉刑、涉民（法院定）、拖欠劳动者报酬（国务院 \hz 省政府）等。
                  \item 强制出境。
                        \begin{enumerate}
                          \item \zd{限期}。个案问题不具有可考性。
                          \item \zd{遣送}。县以上公安 \hz 边检，\zd{1 到 5 年}不得入境。
                          \item \zd{驱逐}。公安部，10 年不得入境。
                        \end{enumerate}
                \end{enumerate}
        \end{enumerate}
  \item 引渡法。
        \begin{enumerate}
          \item 外交保护。在境外的本国人，\zd{属人管辖}的体现。\zd{不以受害者请求为条件}。
                \begin{enumerate}
                  \item 侵害行为要直接 \hz 间接\zd{归因国家}。
                  \item 必须\zd{先}用尽当地救济。
                  \item \zd{国籍继续}原则。外交保护结束前，受害者要持续拥有保护国国籍。
                \end{enumerate}
          \item 引渡。\zd{是国家权利}，不是义务。国家和国家之间，罪犯 \hz 嫌犯。\al{甲国有义务将乙引渡给中国 \cuo}
                \begin{enumerate}
                  \item 前提条件。\bs{或引渡条约或互惠承诺}。\al{中乙没有引渡条约，乙提出引渡请求应当作出互惠的承诺}。
                  \item \zd{应当拒绝}引渡。\bs{国民双重政治犯，审完过期军事犯，酷刑程序缺席判}。军事犯，两边之一构成即可。
                  \item 可以拒绝引渡。具有刑事管辖 + 正在进行刑事诉讼。
                  \item 引渡对外外交部一张嘴。\zd{高院裁，最高法核}，限制追诉最高检（定罪），限制量刑最高法。
                  \item 罪名特定。\zd{什么罪引的什么罪判}。
                  \item 非经引出国同意\zd{不得转引渡}。要看前提。
                        \al{中国公民张某杀害甲国公民后逃亡乙国。如乙将张引渡到中国后，未经乙国同意，我国不得再将张某引渡到甲国 \cuo 张是中国人，拒绝引渡无需乙同意}。
                \end{enumerate}
                \begin{kk}{政治犯}
                  一般由被请求引渡国决定，例外。
                  \begin{enumerate}
                    \item 战争罪。
                    \item 危害人类罪。种族灭绝 \hz 种族隔离。
                    \item 劫机罪。
                    \item 侵害外交代表、受国际保护人员。
                  \end{enumerate}
                  \sjt 不能以政治犯为由拒绝引渡（这些罪伤天害理，无权获得庇护、保护），但\zd{不意味必须引渡}。
                  \tcblower
                  2005 年联合国反腐败公约、2003 年联合国打击跨国国有组织犯罪公约。
                  \begin{enumerate}
                    \item 公约\zd{可以但不必然作为}缔约国之间产生引渡义务的法律依据。
                    \item 有义务接引渡材料，拒绝要说明理由。
                    \item \zd{必须承担义务}。不可以以双边条约对抗在多边公约下承担的义务。
                    \item 没有改变本国国民不引渡原则。\zd{仅以本国国民为由拒绝引渡的，要么引渡 \hz 本国进行起诉}。
                  \end{enumerate}
                \end{kk}
          \item 庇护。\zd{受追害}的境内的外国人，\zd{属地管辖}的体现。\zd{是国家的权利不是义务}。
                \begin{enumerate}
                  \item 给予其居留权 + 拒绝将其引渡。
                  \item 无权庇护 4 种罪同引渡。
                  \item 外交庇护/域外庇护。\zd{非法}（我国不承认），本质是双方对峙精神喊话。
                \end{enumerate}
        \end{enumerate}
\end{enumerate}
